% SOURCE_STATEMENT: "The long-time behavior of the solution of the Cauchy problem for a great many integrable systems on the line is now well-understood, using, for example, the Riemann-Hilbert/steepest descent method. The method depends on full knowledge of the nature of the spectrum of the associated Lax operator. For systems such as KdV, MKdV, defocusing NLS and the Toda lattice, for example, one is able to describe the solution asymptotically in complete detail for general initial data. But for focusing NLS, where the associated AKNS Lax operator $T$ is non-self-adjoint (here we make contact with Problem 7), the situation is different. For generic $T$ (i.e. an open dense set of $T$'s in any reasonable topology) the spectrum consists of the real line, where the spectrum is absolutely continuous, together with a finite number of simple eigenvalues (corresponding to solitons) off the real axis. The analysis of the long-time behavior of focusing NLS with such generic initial data proceeds in a straightforward manner similar to KdV, MKdV, etc. For general initial data, however, the situation is more complicated. For example, let $z_0>0$ be any positive number and let $D$ be any arbitrarily small open disk in the complex plane centered at $z_0$, such that $D\\subset\\{z\\,:\\,\\mathbb{R}e z>0\\}$. Let $D^+$ denote the intersection of $D$ with the upper half plane, and let $u(z)$ be an arbitrary function analytic in $D^+$, and continuous in $\\overline{D^+}$. Let $B=\\{z\\in D^+\\,:\\, u(z)=0\\}$. Then there exists (Zhou) an AKNS operator $T$ with infinitely smooth, rapidly decaying coefficients with the property that each point in $B$ is an $L^2(\\mathbb{R})$-eigenvalue of $T$. In other words, there exist (non-generic) operators $T$ with Schwartz space coefficients which have $L^2$ spectrum accumulating on the real line at an essentially arbitrary rate. It is a very interesting question to determine what effect such singularities would have on the long-time behavior of the solution of NLS. In particular, recalling that focusing NLS provides a model for data transmission along communication cables, are the effects measurable? The difficulty that we encounter here is not limited to situations where the associated Lax operator is non-self-adjoint. Even in situations when the associated operator is self-adjoint, but of order greater than two, similar difficulties can arise. This is true, in particular, for the Boussinesq equation, where the associated Lax operator is third order. The long-time behavior of the solutions of the Boussinesq equation with general initial data is a very interesting problem with many challenges. Even in the case with generic initial data the situation is only partially understood."
% Preparation and final self-audit: CONTRIBUTING.md prompts C and D; source/audit artifacts retained in the local physics_statements_2026_10_11 batch.
% !TeX program = lualatex
\documentclass[11pt]{article}
\usepackage[margin=25mm]{geometry}
\usepackage{fontspec}
\setmainfont{Latin Modern Roman}
\usepackage{amsmath,amssymb,mathtools}
\usepackage{unicode-math}
\setmathfont{Latin Modern Math}
\usepackage{xurl,hyperref}
\hypersetup{colorlinks=true,urlcolor=blue}
\setcounter{secnumdepth}{0}
\setlength{\parindent}{0pt}
\setlength{\parskip}{5pt}
\setlength{\emergencystretch}{3em}
\newcommand{\sref}[2]{\hyperlink{source:#1:#2}{[S#2]}}
\newcommand{\cataloguescope}[1]{\paragraph{Recorded target.}#1}
\newcommand{\R}{\mathbb R}
\newcommand{\C}{\mathbb C}
\newcommand{\N}{\mathbb N}
\newcommand{\Z}{\mathbb Z}
\newcommand{\Q}{\mathbb Q}
\newcommand{\F}{\mathbb F}
\newcommand{\eps}{\varepsilon}
\begin{document}
% UnsolvedMath ID 8000008; source catalogue code AMR-079-0008
\setcounter{section}{8000007}
\section{Long-time behavior with non-generic initial data}\label{8000008}
{\small UnsolvedMath ID 8000008 · Mathematical Physics}
\subsection{Definitions and mathematical statement}

\paragraph{Shared notation.}$\mathbb N=\{1,2,\ldots\}$; the standard number fields have their usual meanings. All additional conventions are stated locally.


\paragraph{Definitions and notation.}For focusing NLS the associated AKNS scattering operator is non-self-adjoint. Discrete eigenvalues off the real axis correspond to soliton data. Accumulation of infinitely many such eigenvalues at the real continuous spectrum creates exceptional scattering singularities. For the Boussinesq equation the Lax operator is third order, and non-generic spectral behavior can likewise obstruct asymptotic analysis.

\paragraph{Catalogue wording retained for traceability.}
The long-time behavior of the solution of the Cauchy problem for a great many integrable systems on the line is now well-understood, using, for example, the Riemann-Hilbert/steepest descent method. The method depends on full knowledge of the nature of the spectrum of the associated Lax operator. For systems such as KdV, MKdV, defocusing NLS and the Toda lattice, for example, one is able to describe the solution asymptotically in complete detail for general initial data. But for focusing NLS, where the associated AKNS Lax operator $T$ is non-self-adjoint (here we make contact with Problem 7), the situation is different. For generic $T$ (i.e. an open dense set of $T$'s in any reasonable topology) the spectrum consists of the real line, where the spectrum is absolutely continuous, together with a finite number of simple eigenvalues (corresponding to solitons) off the real axis. The analysis of the long-time behavior of focusing NLS with such generic initial data proceeds in a straightforward manner similar to KdV, MKdV, etc. For general initial data, however, the situation is more complicated. For example, let $z_0>0$ be any positive number and let $D$ be any arbitrarily small open disk in the complex plane centered at $z_0$, such that $D\subset\{z\,:\,\mathbb{R}e z>0\}$. Let $D^+$ denote the intersection of $D$ with the upper half plane, and let $u(z)$ be an arbitrary function analytic in $D^+$, and continuous in $\overline{D^+}$. Let $B=\{z\in D^+\,:\, u(z)=0\}$. Then there exists (Zhou) an AKNS operator $T$ with infinitely smooth, rapidly decaying coefficients with the property that each point in $B$ is an $L^2(\mathbb{R})$-eigenvalue of $T$. In other words, there exist (non-generic) operators $T$ with Schwartz space coefficients which have $L^2$ spectrum accumulating on the real line at an essentially arbitrary rate. It is a very interesting question to determine what effect such singularities would have on the long-time behavior of the solution of NLS. In particular, recalling that focusing NLS provides a model for data transmission along communication cables, are the effects measurable? The difficulty that we encounter here is not limited to situations where the associated Lax operator is non-self-adjoint. Even in situations when the associated operator is self-adjoint, but of order greater than two, similar difficulties can arise. This is true, in particular, for the Boussinesq equation, where the associated Lax operator is third order. The long-time behavior of the solutions of the Boussinesq equation with general initial data is a very interesting problem with many challenges. Even in the case with generic initial data the situation is only partially understood.
\paragraph{Scope and source qualification.}The generic-data results and exceptional-data examples in the original text remain context. One must cover the specified resonant cases rather than substitute the regular reflection-coefficient assumptions of the established method.
\subsection{Short English statement}
Determine long-time behavior of integrable PDEs for exceptional initial data that produce singular scattering behavior.
\subsection{Sources}
\begin{itemize}
\item[\hypertarget{source:8000008:1}{S1}] ULAM.AI and the UnsolvedMath Contributors, UnsolvedMath, record 8000008 (AMR-079-0008). Catalogue curation under CC BY 4.0; definitions and English exposition added locally.
\url{https://huggingface.co/datasets/ulamai/UnsolvedMath}.
\item[\hypertarget{source:8000008:2}{S2}] Percy Deift, Some open problems in random matrix theory and the theory of integrable systems (2007), arXiv:0712.0849, Problem 8.
\url{https://arxiv.org/abs/0712.0849}.
\end{itemize}
\label{end:8000008}
\end{document}
